\documentclass[10pt]{article}

\usepackage[margin=1in]{geometry}
\usepackage{hyperref}
\usepackage{booktabs}

\title{Vision-Only Latent Diffusion Super-Resolution without T2I Pretraining}
\author{sr-diffusion contributors}
\date{Work-in-progress technical report}

\begin{document}
\maketitle

\begin{abstract}
This report describes a work-in-progress x4 super-resolution project. The goal
is to train a vision-only latent diffusion SR model directly, without importing a
pretrained text-to-image diffusion backbone. The current implementation is in
Stage 1: learning a factor-4 VAE/autoencoder on 512px HR crops so later stages
can train in latent space.
\end{abstract}

\section{Objective}

The target task is x4 restoration from low-resolution inputs:
\[
128 \times 128 \rightarrow 512 \times 512
\]
with a later extension to
\[
192 \times 192 \rightarrow 768 \times 768.
\]
The intended model should support photo and anime/illustration domains in one
model through explicit domain conditioning.

\section{Planned Method}

The planned pipeline is:
\[
\mathrm{HR\ image} \rightarrow \mathrm{VAE} \rightarrow \mathrm{HR\ latent}.
\]
During SR training, a degraded LR image is encoded by a condition encoder. A
conditional diffusion U-Net then denoises HR latents using LR features, timestep
embeddings, and domain embeddings. The VAE decoder maps the final latent back to
the x4 HR output.

\section{Current Implementation}

The repository currently implements the Stage 0--1 foundation:
\begin{itemize}
  \item manifest-based photo/anime dataset loading,
  \item on-the-fly x4 degradation,
  \item factor-4 AutoencoderKL,
  \item PyTorch/ROCm training with bf16 autocast,
  \item W\&B online logging,
  \item fixed validation sample logging for LR, GT, and HR reconstruction,
  \item validation metrics for loss, reconstruction, KL, MSE, and PSNR.
\end{itemize}

\begin{table}[h]
\centering
\begin{tabular}{ll}
\toprule
Item & Current value \\
\midrule
Training data & 10,000 photo images \\
Validation data & 100 photo images \\
Input crop & 512px HR \\
Latent size & 128px, 16 channels \\
Batch size & 16 \\
Max steps & 100,000 \\
Eval cadence & Every 1,000 steps \\
Sample cadence & Every 500 steps \\
\bottomrule
\end{tabular}
\end{table}

\section{Evaluation}

Current evaluation is for the VAE only:
\[
\mathrm{HR} \rightarrow \mathrm{encode} \rightarrow \mathrm{decode}
\rightarrow \mathrm{reconstruction}.
\]
The validation set is fixed, and visual samples are also fixed so that W\&B
comparisons show real progress over time rather than different random crops.

\section{Roadmap}

\begin{enumerate}
  \item Stage 1: finish and select the VAE checkpoint.
  \item Stage 2: train deterministic LR-to-HR-latent pretraining.
  \item Stage 3: train conditional latent diffusion for SR.
  \item Stage 4: add perceptual/GAN fine-tuning if needed.
  \item Stage 5: distill to fewer inference steps.
  \item Stage 6: evaluate checkpoints with fixed A/B preference tests.
\end{enumerate}

\section{Limitations}

This is not yet a complete SR diffusion model. The current training run only
learns the VAE reconstruction model. Anime/illustration data has not yet been
added to the active 10k training set, and no claim is made that this is a novel
published method. This is a practical engineering report for an active
experiment.

\end{document}
